% appx-roadmap.tex — Appendix F: open obligations + roadmap.

\chapter{Open Obligations and Refinement Roadmap}
\label{ch:appx-roadmap}

\noindent\emph{The open analytical refinements that would tighten
the constructive mass-gap bound, the structural extensions to
$\SUthree$, and the future-work items the manuscript names but
defers.}

The manuscript closes the loop on the $\SUtwo$ instance of the
Clay statement (\Cref{ch:mass-gap-discussion},
\Cref{sec:deliverables}): existence (C1) via the projective-limit
construction, non-triviality (C2) by construction, mass gap (C3)
via the explicit $\massgap \ge 999/1000$.  This appendix
catalogues what remains: analytical refinements that would tighten
the bound, structural extensions that would broaden it, and
deliberately deferred items.

\section{Analytical refinements (bound tightening)}
\label{sec:roadmap-tightening}

\marginkey{The factor-of-$3.7$ gap}%

The empirical Teper ratio
$m(0^{++}) / \sqrt{\sigma} \approx 3.74$
\cite{Teper1998Glueballs,LuciniTeperWenger2004} versus the
constructive $\massgap / \sqrt{\sigma_0} \approx 1.0$ gives a
looseness factor $\sim 3.7$.  The named tightening targets:

\begin{itemize}
  \item \textbf{Sharper decay-to-gap relation.}  The OS4 cluster
    rate $\sighinf$ is identified with the spectral gap by the
    linear-decay-to-gap relation $m \ge \sqrt{\sighinf}$.  A
    sharper Brydges--Federbush-style estimate that exploits the
    full propagator structure could in principle give
    $m \ge c \cdot \sighinf$ with $c$ larger than $1$.
    \emph{Estimated tightening:} factor $\sim 2$.

  \item \textbf{Polymer-residual $\cpolymer$.}  The
    triangle-inequality worst case $\cpolymer = 2$ used in
    \Cref{ch:sigma-hat-trajectory}, Theorem~\ref{thm:eta-coupling}
    can be reduced to $\cpolymer = 1$ via a sign-preservation
    lemma on the cluster-expansion residual.
    \emph{Estimated tightening:} halves the per-step $\eta$ and
    improves $\massgap$ to $\approx 0.99955$ (relative to the
    baseline $0.999$).

  \item \textbf{$\Kcap$ union-bound looseness.}  The bound
    $\Kcap \le \Ksf + \Klf$ of \Cref{ch:overlap-kp},
    Theorem~\ref{thm:overlap-glue}, is loose by an empirical
    factor up to $\approx 1.18$ (\cref{conj:K-cap-sharpness}).
    A sharper sheaf-domination-gluing bound would tighten the
    chain.  \emph{Estimated tightening:} factor $\sim 1.18$.

  \item \textbf{Direct $\clambda$ extraction.}  The Casimir-gap
    identification $\clambda = \tfrac34$ from
    \Cref{ch:small-field-kp}, Theorem~\ref{thm:c-lambda-casimir}
    is a project-derived synthesis; a direct extraction of the
    spectral-gap coefficient from the primary Bałaban literature
    against matched conventions could refine the value.
    \emph{Estimated tightening:} indirect, dependent on the
    refinement magnitude.

  \item \textbf{Refined $\muplaq$.}  The default $\muplaq = 16$ is
    literature-shape; a tighter plaquette-polymer-counting bound
    on the 4D lattice would shrink the large-field convergent
    threshold, expanding the convergent regime and (indirectly)
    tightening $\Klf$ at the canonical pilot.
    \emph{Estimated tightening:} indirect.
\end{itemize}

None of these is a blocker for the existence statement; each
would refine the quantitative bound but leave its positivity
intact.

\section{Structural extensions}
\label{sec:roadmap-extensions}

\marginkey{$\SUthree$ and beyond}%

The framework is group-agnostic at the abstract level
(\Cref{ch:mass-gap-discussion}, \Cref{sec:su-n-extension}).  The
extensions:

\begin{itemize}
  \item \textbf{$\SUthree$ pilot.}
    \Cref{conj:su3-pilot} of \Cref{ch:mass-gap-discussion}: there
    exists $g_0$ such that the $\SUthree$ pilot validator
    (the analog of \Cref{ch:pilot-su2}'s) closes at $K = 3$.  The
    work required is mechanical: rerun the Phase-1--3 numerics
    with $\SUthree$ representation-theoretic and heat-kernel
    inputs.  No new analytical open question is expected.

  \item \textbf{General $\SUN$.}  The same framework extends to
    arbitrary compact simple gauge groups; the
    representation-theoretic inputs change but the structural
    arguments are unchanged.

  \item \textbf{Confinement (Clay-statement bonus condition).}
    The Clay statement's bonus condition (area-law Wilson loops
    at \emph{all} sizes for $\SUthree$ with quarks) is not
    addressed.  The current framework delivers area-law in the
    convergent regime but does not extend to strong coupling.
    A separate constructive-confinement argument is the natural
    follow-on.

  \item \textbf{Glueball-spectrum identification.}  The bound
    $\massgap$ is the spectral gap of $\Hop$, not the
    lightest-glueball mass per se.  Spin-and-parity-resolved
    spectroscopy beyond the present scope.
\end{itemize}

\section{Constructive-analysis layer}
\label{sec:roadmap-creal}

\marginkey{$\Qp \to \mathbb{R}$}%

The manuscript works in $\Qp$ throughout; every numerical bound
is a rational with explicit denominator.  Promoting the
$\Qp$-rational bounds to genuine constructive-reals statements
would close several deliberate gaps:

\begin{itemize}
  \item Universality of $\sighinf$: \cref{conj:sigh-universality}
    of \Cref{ch:continuum-limit} expects $\sighinf$ to be
    $g_0$-independent in the asymptotic-freedom limit.  A
    constructive-reals proof would establish this universality
    sharply.
  \item Real-valued $\massgap$: the manuscript's
    $\massgap \in \Qp$ is a strict lower bound on the real-valued
    continuum-limit mass gap.  The two are conjecturally close at
    the canonical pilot.  Closing the gap between them requires
    Cauchy-completion machinery at the
    constructive-reals layer.
  \item OS reconstruction in Lean: the
    \Cref{ch:os-reconstruction}'s
    Hilbert-space construction is paper-level, with Lean
    delivering only the type signature.  A full Lean
    formalization (Stone's theorem, spectral theorem for
    self-adjoint operators) is a major undertaking deferred to
    the parallel constructive-reals layer.
\end{itemize}

\section{Remaining loose ends}
\label{sec:roadmap-loose-ends}

\marginkey{Specific deferrals}%

\begin{itemize}
  \item \textbf{Full $O(4)$ rotation restoration.}
    \Cref{ch:os-axioms} verifies hypercubic $H_4$ invariance for
    $\muinf$; the full $O(4)$ rotation invariance requires a
    separate restoration argument that is the paper-level
    deliverable of a future revision.

  \item \textbf{Two-plaquette correlation coefficient $\zeta$.}
    The chaining factor of \Cref{ch:chaining-estimate} uses
    $\zeta = \tfrac12$ from the leading-order
    SU(2)-heat-kernel two-plaquette joint expansion.  Refining
    $\zeta$ against the full two-plaquette integral would tighten
    the chaining factor and amplify the chaining-tightening
    contribution to $\massgap$.

  \item \textbf{Sign-preservation lemma on cluster-expansion
    residual.}  The $\cpolymer = 1$ tightening of
    \Cref{sec:roadmap-tightening} is conditional on a
    sign-preservation lemma for the SU(2) heat-kernel polymer
    activities under $\Tb$.  No proof yet; sign-preservation is
    not generic for cluster-expansion polymer activities.

  \item \textbf{Tightness of Conjecture~\ref{conj:gaussian-sharpness}.}
    The empirical observation $L_K \le 1.18$ across the canonical
    convergent regime supports the universal bound $L_K \le 2$,
    but a proof of the universal bound is open.
\end{itemize}

\section{Project-completion outlook}
\label{sec:roadmap-completion}

\marginkey{What remains}%

In its current state the manuscript carries an existence proof of
the $\SUtwo$ mass gap, $\massgap \ge 999/1000$ in dimensionless
units, certified by a Lean witness with axiom-gate
\codepath{[propext]} alone.  Going forward, the natural ordering
of refinements:

\begin{enumerate}[label=(R\arabic*)]
  \item Close the $\SUthree$ pilot (\Cref{conj:su3-pilot}).
    Highest priority because it extends the existence statement
    to the physically relevant gauge group; mostly mechanical
    work over the existing framework.
  \item Address the polymer-residual tightening
    ($\cpolymer = 1$): doubles the
    margin on $\massgap$ at the canonical pilot.
  \item Address the decay-to-gap tightening: factor-$\sim 2$
    on $\massgap$.
  \item Build the constructive-reals layer that promotes the
    $\Qp$-rational bounds to real-valued statements; close the
    universality conjectures.
  \item Restore full $O(4)$ rotation invariance in
    \Cref{ch:os-axioms}.
\end{enumerate}

Each is a well-defined refinement task on top of the framework
laid out in this manuscript.  None of them changes the existence
claim: a constructive, strictly positive, $\Qp$-rational
$\SUtwo$ mass gap has been established.
